\section{Introduction}
\label{sec:introduction}
Large language models (LLMs) are increasingly deployed as agents that act on
external software through \emph{tools}. These tools allow models to access
files, execute programs, query databases, issue network requests, and interact
with external services, while orchestration frameworks translate model outputs
into structured tool invocations. This transition substantially increases the
security consequences of model failures: an unsafe output can now become an
operation executed by conventional software with access to sensitive
resources resulting in jailbreaks
~\cite{zou2023universal}, indirect prompt injection
~\cite{greshake2023not}, tool poisoning~\cite{mcp-scanner}, tool
hijacking~\cite{toolhijacker}, and cross-tool attacks
~\cite{les-dissonances}. 


In this paper, we ask a fundamental question: \emph{what level of trust should be established between the LLM and the tool?} We surveyed major orchestration frameworks and protocols by tracing tool registration, argument checking, and dispatch, and found that they largely enforce only structural constraints encoded in schemas, while semantic constraints expressed in natural-language metadata remain unenforced. For example, LangChain and MCP can check that \texttt{path} is a string, but not that a path described as ``within the workspace'' actually remains within that workspace~\cite{langchain-tools,mcp-tools}. This creates a gap between the interface presented to the LLM and the guarantees provided to the tool. The gap is security-critical because LLM-generated arguments may still derive from adversarial prompts, retrieved content, or other tools. As we show in Section~\ref{sec:trust-boundary}, semantic validation is only one of several missing protections across this boundary. We therefore argue that \emph{LLM-generated tool arguments must remain untrusted at execution time, as tool implementations cannot be assumed to enforce the security properties implied by their documented interfaces}.

Existing techniques are poorly suited to answering this question at scale.
Current analyses of LLM tools primarily encode predefined weakness classes,
specific information flows, or particular agent--tool attacks
~\cite{mcp-first-glance,mcp-seclint,mcp-scanner,les-dissonances}. Traditional
fuzzing explores broader input spaces, but generally relies on generic oracles
such as crashes, memory-safety violations, or assertions. Many tool
vulnerabilities produce no such signal. A file tool may successfully access a
resource outside its intended workspace, or a network tool may successfully
contact a prohibited destination. Detecting these failures requires knowing
\emph{what an argument represents} and whether execution respects the
corresponding security property.

Our key insight is that LLM tools already expose the semantic information
needed to construct such tests. Tool metadata exists precisely to describe to
the model what a tool and its arguments mean. Two arguments may both have type
\texttt{string}, while one represents a filesystem path and another a URL.
Their relevant security properties, adversarial inputs, and runtime effects
are correspondingly different. We can therefore repurpose existing tool
metadata as a source of security specifications for testing the implementation
behind the interface.

We present \project{}, a tool-aware framework for automatically discovering
semantic security vulnerabilities in LLM tools. \project{} introduces
\emph{semantics-aware sanitizers}, a reusable corpus that associates semantic
input classes with CWE-guided security properties, adversarial inputs, and the
runtime observations needed to identify violations. Given a registered tool,
\project{} maps its arguments to applicable sanitizers and specializes those
sanitizers using the semantics documented in the tool metadata. An LLM assists
with interpreting natural-language metadata, but is not used as the security
oracle.
\project{} then automatically generates schema-valid fuzzing harnesses and
semantic oracles. When implementation-level visibility is available,
\project{} uses runtime introspection to observe security-sensitive effects
and determine whether they violate the selected sanitizer. For opaque tools,
it generates a reference proxy encoding the expected security behavior and
uses differential checking to identify inconsistencies between the advertised
contract and the target implementation. Thus, rather than asking only whether
an execution crashes, \project{} determines whether a tool behaves securely
according to the semantics of its own interface.

We evaluated \project{} on 569 tools across 11 open-source repositories using
66 semantics-aware sanitizers. \project{} discovered 14 previously unknown
bugs across five repositories, spanning arbitrary file writes, MCP server
crashes, denial of service, incorrect tool descriptions, and other semantic
security failures. Seven bugs are input-independent and require no particular
triggering input. These findings show that securing model behavior and
orchestration alone is insufficient: the software ultimately executing
LLM-generated actions frequently fails to enforce the security properties its
interface advertises. Our contributions are as follows:

\begin{itemize}
    \item We systematically analyze the LLM--orchestrator--tool trust boundary
    and identify several missing protections that allow attacker-influenced
    values to reach privileged tool implementations, establishing why
    LLM-generated arguments must remain untrusted at execution time.

    \item We introduce \emph{semantics-aware sanitizers} and design
    \project{}, which derives tool-specific security properties from existing
    metadata, generates schema-valid fuzzing harnesses, and detects violations
    using runtime introspection and black-box differential checking.

    \item We evaluate \project{} on 569 tools across 11 repositories and
    discover 14 previously unknown bugs across five repositories,
    demonstrating that real-world LLM tools frequently fail to enforce the
    semantic contracts exposed by their interfaces.
\end{itemize}

